\section{Limitations}
\label{sec:limitations}

We describe the design's limitations honestly and in one place so that a
reader can weigh what the paper does and does not claim.

\subsection{Scale and Deployment}
\label{sec:limits:scalability}

The current prototype is single-node and in-memory. All measurements are at
scales of $100\,000$ facts or below; behaviour at larger scales is
extrapolated from the complexity analysis of Section~\ref{sec:complexity}
and not directly measured. Global-section computation in the worst case is
$O(|\mathcal{C}| \cdot N^2)$, which imposes a natural ceiling on the size
of context poset that the system can support on a single node. This
ceiling is compounded by a concretely measured one: Section~\ref{sec:eval:memory}
reports \rangememkg{} higher resident memory per fact than the KG
baseline, so the single-node memory budget is consumed faster by SFDB than
by an equivalent reified store, independent of the gluing worst case.
Distributing the store across nodes would require a partitioning strategy
for the context poset and a distributed restriction-map protocol; both are
outside the scope of this paper.

\subsection{Mathematical Assumptions}
\label{sec:limits:math}

The formal model assumes a finite context poset equipped with the
Alexandrov topology. This is the right setting for the intended workload,
but it is also the setting in which the sheaf condition holds vacuously.
Nothing in the paper claims a non-trivial sheaf theorem. The gluing
operation in the implementation is a defensive consistency check rather
than a step in a proof; in the presence of genuinely contradictory local
data (which our schema currently forbids), the algorithm would report the
conflict rather than resolve it.

A further structural assumption deserves to be stated explicitly: the
dot-separated path encoding makes the context hierarchy a \emph{tree}
(each context has exactly one parent), so incomparable contexts never
share a descendant and meets of incomparable contexts do not exist
(Section~\ref{sec:formal_definitions}). Real-world context structures
are not always tree-shaped --- a policy that holds in the intersection
of two jurisdictions, or an experimental condition scoped by both a
protocol and a lab, is naturally a DAG node with two parents. Supporting
DAG-shaped posets would preserve the Alexandrov construction (which
needs only a preorder) but would invalidate the prefix-string encoding,
the longest-common-prefix join, and the single-ancestor-chain
$O(d)$ insert path, so we scope it out here and note it as a design
extension rather than pretending the tree assumption is costless.

\subsection{Implementation Limitations}
\label{sec:limits:implementation}

The prototype does not currently support concurrent transactions, on-disk
persistence beyond the built-in SQLite backing store, distributed
execution, streaming or incremental fact ingestion, or SPARQL compliance.
An earlier draft of this paper described the query language as "a small
subset of SPARQL"; that overstated it. It is a bespoke query DSL whose
surface syntax borrows SPARQL's \texttt{SELECT}/\texttt{WHERE} keywords and
\texttt{?var} notation but does not parse SPARQL's graph-pattern braces or
its \texttt{FROM} clause, together with a few sheaf-native operations
(restriction, stalk lookup). The typed \texttt{Query} object used
internally by both engines and by the benchmark harness
(Section~\ref{sec:evaluation}) does not depend on this text surface. A
real SPARQL front end and a proper transactional layer are engineering
work that would be needed for any real deployment but that would not
change the paper's structural findings.

\subsection{Baseline and Benchmark Limitations}
\label{sec:limits:benchmark}

The KG baseline is a research prototype implemented by the same team as
the sheaf engine. It is a controlled reference point, not a
state-of-the-art RDF store. Both engines share the same serialisation
and canonicalisation layers, and both persist inserts to in-memory
SQLite --- but the \emph{query} paths differ in more than indexing
strategy: the KG baseline answers queries through SQLite, while SFDB
answers them from in-process Python indexes and touches SQLite only on
the write path. A comparison of the two alone therefore conflates two
variables --- context-indexed storage versus reification, and dict
indexes versus SQLite round trips. The evaluation separates them with
the KG-mem control (Section~\ref{sec:eval:control}): a variant of the
KG baseline with identical reification and query logic whose dictionary
tables and triple indexes are Python dicts, so that the SFDB-vs-KG-mem
ratios isolate what context-indexed storage buys with the storage-layer
difference removed.

Both the KG baseline and the KG-mem control are, however, still
research prototypes by this project's authors, and that ceiling has
since been addressed directly rather than left as a caveat: three real
production stores, Apache Jena TDB2, OpenLink Virtuoso, and Ontotext
GraphDB, are each
compared on all five query classes at every
scale (Section~\ref{sec:eval:jena}, Section~\ref{sec:eval:virtuoso},
Section~\ref{sec:eval:graphdb}),
which is the natural next step this
section previously called out as future work. The result is a genuine
crossover, reproduced independently against all three --- SFDB faster on
LOOKUP/GLOBAL/CONTEXT, the production store faster on
TEMPORAL from $10^4$ facts --- not a clean win either way, which
we take as more credible than a one-sided result would have been, and
the agreement between three unrelated optimisers as evidence the pattern
generalises well beyond one vendor. What
remains open is breadth, not existence: three stores is not an
exhaustive survey, and LUBM~\cite{guo2005lubm} (Section~\ref{sec:eval:lubm}) has
since been run both against the three in-house engines and against all
three external stores at once (Section~\ref{sec:eval:lubm-external}),
reproducing the identical pattern against every one; BSBM~\cite{bizer2009bsbm} and
WatDiv~\cite{aluc2014watdiv} have not been tried at all, nor has LUBM at
the much larger scales published LUBM results use.
The dataset-bias half of this concern has since been addressed
directly rather than left as a caveat, and the answer is not the tidy
one the LOOKUP/GLOBAL results alone would suggest: a real-world case
study (Section~\ref{sec:eval:wikidata}, \wikidatanumfacts{} genuine
Wikidata facts) and the LUBM comparison both
reproduce the LOOKUP, GLOBAL, and CONTEXT advantage, while TEMPORAL
lands somewhere different on each --- beating both baselines on
Wikidata, roughly at parity with them on LUBM's two largest scales,
neither matching the synthetic benchmark's own result. The synthetic datasets, while
parameterised, do not exhibit all pathologies of organic knowledge
graph data, and these two comparisons are direct evidence of exactly
that for TEMPORAL specifically --- and, for CONTEXT, direct evidence of
something else: LUBM's topology exposed a genuine performance bug in
SFDB's own query resolution (Section~\ref{sec:discussion:threats}'s
tenth finding) that neither the synthetic generator nor Wikidata's data
shape had triggered, which is itself an argument for dataset diversity
in an evaluation independent of what any single comparison's headline
result turns out to be.

A further limitation of the KG-mem control deserves to be named rather
than left implicit: KG-mem removes the SQLite round trip, but it still
\emph{reconstructs} each fact from its decomposed triples on every query
(\texttt{\_reconstruct\_fact} in the KG engine, unchanged by the
storage-layer swap) rather than returning a stored whole-fact object the
way SFDB's stalk index does. Nothing in principle stops a reified store
from also maintaining a fact-id-to-materialised-fact cache alongside its
triple indexes --- which is architecturally close to what SFDB's stalk
index is --- at which point some of the LOOKUP, GLOBAL, and CONTEXT advantage that
survives the KG-mem control in Section~\ref{sec:eval:control} would
likely narrow further, for the same memory cost SFDB already pays. This
concern is somewhat allayed by the Jena, Virtuoso, and GraphDB results:
all three are
exactly such mature, materialising, professionally engineered indexes,
independently built by three different organisations, and all three
still lose to SFDB on LOOKUP, GLOBAL, and CONTEXT at every scale measured despite
also paying HTTP overhead SFDB does not --- so the advantage on those three
classes does not appear to be merely an artefact of comparing against an
under-optimised control, or of one vendor's specific engineering
choices. The control isolates the SQLite-versus-dict variable cleanly; it does not
isolate reconstruction-from-parts versus retrieval-of-a-whole, which is
a second, unremoved variable the LOOKUP, GLOBAL, and CONTEXT numbers still
partially reflect. (CONTEXT's advantage is dominated by candidate-set
lookup rather than reconstruction, so this caveat bears on it more weakly
than on LOOKUP and GLOBAL, but every result row is still reconstructed
once found.) We do not add a second control for this in the
current evaluation and name it here as an open question rather than
letting the existing control's precision imply more than it isolates.
